% Q1-Q2 问题分析；经 ACCT 于 2026-09-11 审核确认。
\section{问题分析}

\subsection{问题一的分析}

问题一要求在已知单日各时段电价、负荷预测功率和光伏预测功率的条件下，确定外网购电与储能充放电计划。首先，需严格保留附件原始行序，将功率统一换算为每个 10 分钟时段的电量；其次，以负荷、光伏、外网和储能建立逐时能量平衡，并用储电量递推关系刻画充放电损耗，同时满足充放电功率、储能容量及日初日末储电量相等等约束。由于购电费用和上述约束均可表示为线性形式，因此可建立以全天计划购电费用最小为目标的线性规划模型，求得一个最优购电与储能调度方案，再按题目指定时段和四小时区间汇总结果。

\subsection{问题二的分析}

问题二在问题一的基础上进一步引入全年负荷与光伏预测偏差，问题的关键由单日确定性调度转变为受信息时序约束的跨日联动决策。一方面，每日 0:00 制定的普通购电计划在日内不能追加，计划不足产生的供电缺口只能以 5 倍电价紧急购电；另一方面，储能实际日末状态将传递至次日，当前决策会影响后续运行。因此，需要在不使用未来实际信息的前提下，同时控制计划不足造成的紧急购电费用和计划过量造成的无效支出。对此，先根据历史负荷—光伏联合预测残差构造多情景日前计划模型，再在日内逐时观测实际数据，采用滚动优化修正普通取电和储能充放电，仅执行当前时段决策，并将实际储电量连续传递至下一日，从而形成日前计划与日内执行相衔接的闭环策略。
